% !TEX root =./ECE444_Practice_main.tex
%
% L16 practice set (The Array Factor and Pattern Multiplication) -- builds the
% array factor from element phasors, reads the uniform closed form, applies
% pattern multiplication, and sizes element spacing against grating lobes.

\fullwidth{\textbf{Learning Objective 3.2: } Derive the array factor of a linear
array by summing element phasors, reduce the uniform case to its closed form,
apply pattern multiplication, and relate element spacing to the visible region
and the onset of grating lobes.
\\[4pt]\normalfont\small Show your work. A \textbf{Documentation} line at the end
is required (write \textbf{None} if you did not collaborate).}

\begin{questions}

%========================================================== Q1
\question \textbf{Building the sum.} Three isotropic elements sit on a line at
$z = 0$, $z = d$, and $z = 2d$ with $d = \lambda/2$. They are fed in phase with
amplitude weights $a_0 = 1$, $a_1 = 2$, $a_2 = 1$. Scan angle $\theta$ is
measured from broadside.

\begin{parts}

	\part Pattern multiplication lets you write the array pattern as one element
	pattern times one sum. State the two conditions the elements must satisfy for
	that factoring to be legal, and name the step in the derivation where each one
	is used.
		\begin{solution}[1.0in]
		The elements must be \textbf{identical} and \textbf{identically oriented}.
		Both are used at the superposition step: only if every element carries the
		same $E_0(\theta)$ can that factor be pulled outside the sum. If the elements
		differ, or one is rotated, each term keeps its own pattern.
		\end{solution}

	\begin{minipage}[t]{0.58\linewidth}
		\part Write the array factor as a function of $\psi = kd\sinp{\theta}$, then
		find the extra path length and the extra phase that element 2 contributes
		relative to element 0 when $\theta = 30\mydeg$.
		\begin{solution}[1.5in]
		\eq{ AF(\psi) &= 1 + 2e^{j\psi} + e^{j2\psi} }
		The extra path is the projection of the element position on the look
		direction:
		\eq{ \Delta r_2 &= 2d\sinp{30\mydeg} = 2(0.5\lambda)(0.5) = 0.5\lambda }
		\eq{ \Delta\phi_2 &= k\,\Delta r_2 = \frac{2\pi}{\lambda}(0.5\lambda) = \pi = 180\mydeg }
		\end{solution}
	\end{minipage}\hfill%
	\begin{minipage}[t]{0.37\linewidth}
		\raggedleft \vspace{12pt}
		$\Delta r_2$ = \ansbox[1.3in]{$0.5\lambda$}
		\\[10pt]
		$\Delta\phi_2$ = \ansbox[1.3in]{$180\mydeg$}
	\end{minipage}

	\begin{minipage}[t]{0.58\linewidth}
		\part Evaluate $\vert AF\vert$ at broadside and at $\theta = 30\mydeg$, and
		give the $30\mydeg$ value in dB relative to the broadside peak.
		\begin{solution}[1.3in]
		Factor the sum: $AF = e^{j\psi}(2 + 2\cos\psi)$, so
		$\vert AF\vert = 2 + 2\cos\psi$.
		\eq{ \theta = 0: \quad \psi &= 0, \qquad \vert AF\vert = 4 }
		\eq{ \theta = 30\mydeg: \quad \psi &= \pi\sinp{30\mydeg} = 90\mydeg }
		\eq{ \vert AF\vert &= 2 + 2\cos(90\mydeg) = 2 }
		\eq{ 20\mylog{2/4} &= -6.0\ \db }
		\end{solution}
	\end{minipage}\hfill%
	\begin{minipage}[t]{0.37\linewidth}
		\raggedleft \vspace{12pt}
		$\vert AF\vert_{\theta=0}$ = \ansbox[1.1in]{$4$}
		\\[10pt]
		Level at $30\mydeg$ = \ansbox[1.1in]{$-6.0\ \db$}
	\end{minipage}

	\part The derivation keeps the phase difference between elements but throws
	away the amplitude difference. Explain in one or two sentences why that is
	consistent, given a $112\ \text{mm}$ array observed at a range of $30\ \m$.
		\begin{solution}[1.0in]
		The path lengths differ by at most the array length, $112\ \text{mm}$, which
		against a $30\ \m$ range changes $1/r$ by about $0.4\%$. That same
		$112\ \text{mm}$ is nearly four wavelengths at X-band, so it changes the phase
		by up to $1400\mydeg$. The pattern is built entirely from phase.
		\end{solution}

\end{parts}

\newpage
%========================================================== Q2
\question \textbf{The course array.} The PHASER carries $N = 8$ patch elements on
a line with spacing $d = 14\ \text{mm}$, uniformly excited and unsteered. Work at
$f = 10.3\ \ghz$.

\begin{parts}

	\begin{minipage}[t]{0.58\linewidth}
		\part Find $\lambda$, the spacing in wavelengths, and the range of the
		array-factor argument $\psi$ over the visible region.
		\begin{solution}[1.5in]
		\eq{ \lambda &= \frac{3\times 10^{8}}{10.3\times 10^{9}} = 29.1\ \text{mm} }
		\eq{ \frac{d}{\lambda} &= \frac{14}{29.1} = 0.481 }
		At broadside $\psi = kd\sinp{\theta}$, and $\sinp{\theta}$ runs over $\pm 1$:
		\eq{ \psi_{\max} &= kd = 2\pi(0.481) = 3.02\ \text{rad} = 173\mydeg }
		so $-173\mydeg \le \psi \le +173\mydeg$, just short of one full period.
		\end{solution}
	\end{minipage}\hfill%
	\begin{minipage}[t]{0.37\linewidth}
		\raggedleft \vspace{12pt}
		$d/\lambda$ = \ansbox[1.1in]{$0.481$}
		\\[10pt]
		$\psi$ range = \ansbox[1.1in]{$\pm 173\mydeg$}
	\end{minipage}

	\begin{minipage}[t]{0.58\linewidth}
		\part Locate every null of the array factor that falls in the visible
		region.
		\begin{solution}[1.4in]
		\eq{ \sinp{\theta_m} &= \frac{m\lambda}{Nd} = \frac{m}{8(0.481)} = \frac{m}{3.85} }
		\eq{ m = 1: \ \theta &= \arcsin(0.260) = 15.1\mydeg }
		\eq{ m = 2: \ \theta &= \arcsin(0.520) = 31.3\mydeg }
		\eq{ m = 3: \ \theta &= \arcsin(0.780) = 51.2\mydeg }
		$m = 4$ needs $\sinp{\theta} = 1.04$, which is not a real angle. Nulls occur
		at $\pm 15.1\mydeg$, $\pm 31.3\mydeg$, $\pm 51.2\mydeg$.
		\end{solution}
	\end{minipage}\hfill%
	\begin{minipage}[t]{0.37\linewidth}
		\raggedleft \vspace{12pt}
		Nulls = \ansbox[1.5in]{$\pm 15.1\mydeg,\ \pm 31.3\mydeg,\ \pm 51.2\mydeg$}
	\end{minipage}

	\begin{minipage}[t]{0.58\linewidth}
		\part Find the first-null beamwidth and the half-power beamwidth, and state
		the first sidelobe level you expect.
		\begin{solution}[1.5in]
		\eq{ \text{FNBW} &= 2\arcsin\lp\frac{\lambda}{Nd}\rp = 2(15.1\mydeg) = 30.1\mydeg }
		\eq{ \text{HPBW} &\approx \frac{0.886\lambda}{Nd} = \frac{0.886}{3.85} = 0.230\ \text{rad} = 13.2\mydeg }
		Uniform excitation puts the first sidelobe about $13\ \db$ down; the exact
		value for $N = 8$ is $-12.8\ \db$, at $\pm 21.9\mydeg$.
		\end{solution}
	\end{minipage}\hfill%
	\begin{minipage}[t]{0.37\linewidth}
		\raggedleft \vspace{12pt}
		FNBW = \ansbox[1.1in]{$30.1\mydeg$}
		\\[10pt]
		HPBW = \ansbox[1.1in]{$13.2\mydeg$}
		\\[10pt]
		SLL = \ansbox[1.1in]{$-12.8\ \db$}
	\end{minipage}

	\part A uniform array of $N$ elements has $N-2$ sidelobes in one period of
	$\psi$. Count the sidelobes that are actually visible for this array, and
	explain the count using your answer to part (a).
		\begin{solution}[1.2in]
		With $N = 8$ there are six sidelobes per period, three on each side of the
		main lobe. The visible region spans $\pm 173\mydeg$ of $\psi$, which is
		$96\%$ of the $\pm 180\mydeg$ period, so all six appear --- the three nulls
		found in part (b) bracket three sidelobes on each side. Nothing is cut off
		because the spacing is just under half a wavelength.
		\end{solution}

\end{parts}

\newpage
%========================================================== Q3
\question \textbf{Pattern multiplication.} Four short dipoles lie end to end
along the array axis, spaced $d = \lambda/2$ and fed in phase. In scan angle the
short dipole's element factor is $EF(\theta) = \cosp{\theta}$. The four-element
uniform array factor has its first sidelobe at $\theta = 47.1\mydeg$, where
$AF_4 = -11.3\ \db$.

\begin{parts}

	\begin{minipage}[t]{0.58\linewidth}
		\part Evaluate the element factor at the sidelobe angle in dB, and find the
		first sidelobe level of the total pattern.
		\begin{solution}[1.4in]
		\eq{ EF(47.1\mydeg) &= \cosp{47.1\mydeg} = 0.681 }
		\eq{ 20\mylog{0.681} &= -3.3\ \db }
		Pattern multiplication is a product, so in dB the two levels add:
		\eq{ \text{SLL} &= -11.3\ \db + (-3.3\ \db) = -14.6\ \db }
		\end{solution}
	\end{minipage}\hfill%
	\begin{minipage}[t]{0.37\linewidth}
		\raggedleft \vspace{12pt}
		$EF$ at $47.1\mydeg$ = \ansbox[1.1in]{$-3.3\ \db$}
		\\[10pt]
		Total SLL = \ansbox[1.1in]{$-14.6\ \db$}
	\end{minipage}

	\begin{minipage}[t]{0.58\linewidth}
		\part The array factor of this array has nulls at $\pm 30\mydeg$ and
		$\pm 90\mydeg$. Which nulls survive in the total pattern, and what is the
		total pattern at $\theta = 90\mydeg$?
		\begin{solution}[1.3in]
		All of them survive: a zero in either factor is a zero in the product. At
		$\theta = 90\mydeg$ both factors vanish, since $\cosp{90\mydeg} = 0$ and the
		$m = 2$ array-factor null also sits there, so the total is zero (a deeper,
		wider null than either factor produces alone).
		\end{solution}
	\end{minipage}\hfill%
	\begin{minipage}[t]{0.37\linewidth}
		\raggedleft \vspace{12pt}
		$\vert F(90\mydeg)\vert$ = \ansbox[1.1in]{$0$}
	\end{minipage}

	\part The element factor changed the sidelobe by $3.3\ \db$ but changed the
	half-power beamwidth by almost nothing. Explain why, in terms of the shape of
	$\cosp{\theta}$ near broadside.
		\begin{solution}[1.2in]
		$\cosp{\theta}$ is flat near $\theta = 0$: at the $\pm 6.3\mydeg$ half-power
		edges of the four-element beam it is still $0.994$, or $-0.05\ \db$, so it
		scales both the peak and the half-power points by the same amount and leaves
		the width where it was. Out at $47\mydeg$ the cosine has fallen by
		$3.3\ \db$, which is why the sidelobe moved. A directive element reshapes the
		skirts of a pattern; the main lobe belongs to the array factor.
		\end{solution}

	\part A colleague proposes fixing a beam that is too wide by replacing the
	short dipoles with more directive elements, keeping $N$ and $d$ the same. Say
	whether that works and what it would actually change.
		\begin{solution}[1.2in]
		It does not fix the beamwidth. The main lobe is set by the array factor,
		whose width follows the total length $Nd$; a more directive element is nearly
		constant across that narrow main lobe and only scales it. What the change
		does deliver is lower sidelobes and less wide-angle radiation. To narrow the
		beam you must increase $Nd$, by adding elements or increasing the spacing
		within the grating-lobe limit.
		\end{solution}

\end{parts}

\newpage
%========================================================== Q4
\question \textbf{Spacing, the visible region, and grating lobes.} A uniform
linear array is steered to $\theta_0$ by a progressive phase, so
$\psi = kd(\sinp{\theta} - \sinp{\theta_0})$ and grating lobes appear wherever
$\sinp{\theta_g} = \sinp{\theta_0} \pm m\lambda/d$.

\begin{parts}

	\begin{minipage}[t]{0.58\linewidth}
		\part An unsteered array uses $d = \lambda$. Find every grating lobe in the
		visible region, and give the width of the visible region in $\psi$.
		\begin{solution}[1.4in]
		\eq{ \sinp{\theta_g} &= 0 \pm \frac{\lambda}{\lambda} = \pm 1 \quad\Longrightarrow\quad \theta_g = \pm 90\mydeg }
		The visible region spans $2kd = 4\pi(d/\lambda) = 4\pi$, which is two full
		periods of the array factor, so exactly one repeat of the main lobe reaches
		each edge of real space.
		\end{solution}
	\end{minipage}\hfill%
	\begin{minipage}[t]{0.37\linewidth}
		\raggedleft \vspace{12pt}
		$\theta_g$ = \ansbox[1.1in]{$\pm 90\mydeg$}
		\\[10pt]
		$\psi$ width = \ansbox[1.1in]{$720\mydeg$}
	\end{minipage}

	\begin{minipage}[t]{0.58\linewidth}
		\part A second array uses $d = 0.75\lambda$ and is steered to
		$\theta_0 = 30\mydeg$. Apply the avoidance criterion, then find any grating
		lobe that is present.
		\begin{solution}[1.5in]
		\eq{ d_{\max} &= \frac{\lambda}{1 + \vert\sinp{30\mydeg}\vert} = \frac{\lambda}{1.5} = 0.667\lambda }
		$0.75\lambda > 0.667\lambda$, so the criterion is violated and a grating lobe
		is in view.
		\eq{ \sinp{\theta_g} &= 0.5 - \frac{1}{0.75} = 0.5 - 1.333 = -0.833 }
		\eq{ \theta_g &= -56.4\mydeg }
		The $+$ branch gives $1.833$, which is not a real angle.
		\end{solution}
	\end{minipage}\hfill%
	\begin{minipage}[t]{0.37\linewidth}
		\raggedleft \vspace{12pt}
		$d_{\max}$ = \ansbox[1.1in]{$0.667\lambda$}
		\\[10pt]
		$\theta_g$ = \ansbox[1.1in]{$-56.4\mydeg$}
	\end{minipage}

	\begin{minipage}[t]{0.58\linewidth}
		\part You must scan to $\pm 45\mydeg$ at $10.3\ \ghz$ with no grating lobe.
		Find the largest usable element spacing in millimeters, and state the
		trade-off that stops you from simply making $d$ much smaller.
		\begin{solution}[1.5in]
		\eq{ d &< \frac{\lambda}{1 + \sinp{45\mydeg}} = \frac{29.1\ \text{mm}}{1.707} = 17.1\ \text{mm} }
		Smaller spacing is safe from grating lobes but shortens the array for a given
		element count, so the beam widens as $0.886\lambda/Nd$; holding the beamwidth
		then requires more elements, more phase shifters, and more channels. Close
		spacing also increases mutual coupling between elements. The usual
		compromise is $0.4\lambda$ to $0.5\lambda$.
		\end{solution}
	\end{minipage}\hfill%
	\begin{minipage}[t]{0.37\linewidth}
		\raggedleft \vspace{12pt}
		$d_{\max}$ = \ansbox[1.3in]{$17.1\ \text{mm}$}
	\end{minipage}

	\begin{minipage}[t]{0.58\linewidth}
		\part In the lab you will turn on every third element of the PHASER, making
		the effective spacing $42\ \text{mm}$ at $10.3\ \ghz$. Predict the grating
		lobe angles for the unsteered array.
		\begin{solution}[1.4in]
		\eq{ \frac{d}{\lambda} &= \frac{42}{29.1} = 1.44 }
		\eq{ \sinp{\theta_g} &= \pm\frac{1}{1.44} = \pm 0.693 \quad\Longrightarrow\quad \theta_g = \pm 43.9\mydeg }
		$m = 2$ would need $\sinp{\theta} = 1.39$, so only the first pair appears.
		Expect three lobes of equal height: the main lobe at $0\mydeg$ and grating
		lobes near $\pm 44\mydeg$.
		\end{solution}
	\end{minipage}\hfill%
	\begin{minipage}[t]{0.37\linewidth}
		\raggedleft \vspace{12pt}
		$\theta_g$ = \ansbox[1.3in]{$\pm 43.9\mydeg$}
	\end{minipage}

\end{parts}

\newpage
%========================================================== Q5
\question \textbf{The array as a sampled aperture.} Two unsteered uniform arrays
are proposed for the same job. Array A has $N = 8$ elements at $d = \lambda/2$.
Array B has $N = 4$ elements at $d = \lambda$.

\begin{parts}

	\begin{minipage}[t]{0.58\linewidth}
		\part Find the total length and the half-power beamwidth of each array.
		\begin{solution}[1.5in]
		\eq{ \text{A}: \ Nd &= 8(0.5\lambda) = 4\lambda }
		\eq{ \text{B}: \ Nd &= 4(1.0\lambda) = 4\lambda }
		\eq{ \text{HPBW} &\approx \frac{0.886\lambda}{4\lambda} = 0.222\ \text{rad} = 12.7\mydeg }
		Both arrays are the same length, so both have the same main lobe.
		\end{solution}
	\end{minipage}\hfill%
	\begin{minipage}[t]{0.37\linewidth}
		\raggedleft \vspace{12pt}
		$Nd$ (both) = \ansbox[1.1in]{$4\lambda$}
		\\[10pt]
		HPBW (both) = \ansbox[1.1in]{$12.7\mydeg$}
	\end{minipage}

	\begin{minipage}[t]{0.58\linewidth}
		\part Find the first sidelobe level and the grating lobe angles of each
		array.
		\begin{solution}[1.5in]
		Both are uniformly excited, so both show the uniform first sidelobe near
		$-13\ \db$ ($-12.8\ \db$ for A at $N = 8$, $-11.3\ \db$ for B at $N = 4$).
		\eq{ \text{A}: \ \sinp{\theta_g} &= \pm\frac{\lambda}{0.5\lambda} = \pm 2 \quad \text{(not real, none)} }
		\eq{ \text{B}: \ \sinp{\theta_g} &= \pm\frac{\lambda}{\lambda} = \pm 1 \quad\Longrightarrow\quad \pm 90\mydeg }
		\end{solution}
	\end{minipage}\hfill%
	\begin{minipage}[t]{0.37\linewidth}
		\raggedleft \vspace{12pt}
		A: $\theta_g$ = \ansbox[1.1in]{none}
		\\[10pt]
		B: $\theta_g$ = \ansbox[1.1in]{$\pm 90\mydeg$}
	\end{minipage}

	\part Array B uses half the elements for the same beamwidth. State which array
	you would build and why, and say what happens to Array B's pattern the moment
	it is steered off broadside.
		\begin{solution}[1.4in]
		Build Array A. Both have the same main lobe because both have the same
		length $Nd$, but Array B radiates a full-strength grating lobe at
		$\pm 90\mydeg$ at broadside, and steering slides that lobe into the forward
		half-space while the intended beam moves the other way --- so power and
		receive sensitivity go where they were not commanded. Saving four channels leaves the array
		unable to steer at all. Element count is what provides grating-lobe
		headroom; length is what sets the beamwidth.
		\end{solution}

\end{parts}

\vspace{1.5em}
\noindent\textbf{Documentation:} \rule[-2pt]{0.6\linewidth}{0.4pt}

\end{questions}
